\section{Mathematical setting and the field taxonomy}\label{sec:math}

The size of a field, the order of the whole curve group, and the order of
the subgroup containing the logarithm must remain separate. Throughout,
$q$ denotes a field cardinality, $r$ a prime subgroup order, $h$ a
cofactor, and $Q=[d]G$ a target in $\langle G\rangle$. A curve over an
83-bit field can contain a much smaller prime subgroup. Its ECDLP is
priced in that subgroup, not in the field size.

\subsection{Generic reference algorithms}

Pollard rho obtains a collision among combinations $[a]G+[b]Q$ and solves
the resulting scalar congruence. The familiar random-collision
normalisation is approximately $\sqrt{\pi r/2}$ group steps under idealised
assumptions. This is a model constant, not a universal implementation
bound. Negation, efficiently computable automorphisms, fruitless cycles,
distinguished-point density and restart policy change observed constants.
Parallel collision search changes throughput and coordination costs
without by itself changing the square-root exponent.\cite{pollard1978,vanoorschot1999}

Baby-step giant-step spends square-root memory to obtain another generic
tradeoff. Kangaroo is appropriate when the unknown scalar is known to lie
in an interval; its scale is the interval width, not automatically $r$.
Pohlig--Hellman decomposes composite orders, so a ``large curve'' whose
largest prime factor is small is not a large prime-order ECDLP instance.
These distinctions are implemented in the native toolkit and in the
challenge infrastructure.\src{CA-README}\src{COST}

Generic-group lower bounds establish asymptotic limitations under an
explicit oracle model. They should not be read as a sharp leading-constant
claim of exactly $\sqrt{r/2}$ operations for every implementation with free
negation. The cost challenge is valuable as an executable metering
instrument; this retrospective does not adopt every explanatory constant
in its README as a separately proved lower bound.

\subsection{Index calculus as a complete method}

Choose a factor base $\B=\{F_1,\ldots,F_B\}$, or a quotient of it under
known symmetries. A relation expresses a sampled point as a short sum of
base points. For example,
\[
 [a]G+[b]Q=\sum_i c_iF_i
 \quad\Longrightarrow\quad
 a+bd\equiv\sum_i c_i\log_G(F_i)\pmod r.
\]
After enough independent relations, linear algebra and a target descent
recover the logarithm. Relation count, matrix rank, and recovered scalars
are separate checks. A large number of valid but dependent rows can leave
the system unable to determine the requested target.

For a random-set heuristic with $B$ points and $m$ unordered summands, the
expected number of decompositions is roughly $B^m/(m!r)$. This is not a
universal yield theorem: repeated points, signs, orbit stabilisers,
cofactors, projection into $\langle G\rangle$ and collisions can change the
count. The programme repeatedly found that exact subgroup and base
conventions matter.\src{BOUNDARY}\src{HEUR}

Semaev's summation polynomials encode the existence of point sums through
their abscissae. They provide algebraic formulations, but replacing a group
problem by a polynomial system does not make the system cheap. The cost
of constructing, solving, lifting and verifying that system is central to
the programme.\cite{semaev2004}

\begin{table}[htbp]\small
\caption{Field and problem regimes in the portfolio. Evidence in one row
does not automatically transfer to another.}
\begin{tabularx}{\linewidth}{@{}p{2.7cm}YY@{}}\toprule
Regime & Structure used & Principal limitation in this record \\\midrule
$E(\F_p)$ & coordinate subsets, multiplicative towers, low-degree maps,
special endomorphisms & no general sub-rho prime-field IC demonstrated \\
Ordinary $E(\F_{2^n})$ & $\F_2$ subspaces, Weil restriction, Boolean
algebra & decomposition degree and total yield remain limiting \\
Koblitz $E/\F_2$ over $\F_{2^n}$ & Frobenius orbits and known subgroup
eigenvalues & rho receives the same available symmetry; setup is material \\
Composite-degree binary fields & proper subfields and GHS descent &
weak classes, genus and explicit maps determine applicability \\
Prime-degree binary extensions & no proper intermediate subfield &
composite-degree GHS conclusions require a new construction \\
$E(\F_{p^k})$, small fixed $k$ & base-field abscissae, Gaudry/Joux--Vitse,
torsion invariants & expensive algebraic constants and linear algebra \\
Higher-genus Jacobians & divisor arithmetic, covers, large primes &
source-subgroup transport and same-problem reference must be verified \\
Supersingular isogeny problems & graph navigation, quaternion and Kani
constructions & an isogeny problem or KEM is not the ordinary ECDLP \\\bottomrule
\end{tabularx}
\end{table}

\subsection{Prime, binary, and extension fields}

In a prime field, Frobenius exists but is the identity on field elements;
there is no proper intermediate subfield or nontrivial $\F_p$-linear
coordinate subspace. Saying that ``$\F_p$ has no Frobenius'' is therefore
an imprecise shorthand. The actual obstacle is the absence of a useful
nontrivial subfield decomposition of the kind available over extensions.
This does not rule out nonlinear sparse sets or other non-generic
algorithms.

In $\F_{2^n}$, an $\F_2$-subspace of dimension $\ell$ supplies a compact
membership condition for a factor base. Weil restriction turns an
extension-field polynomial into Boolean equations. The behaviour of the
resulting ideal depends on the formulation and product spaces, not merely
the nominal field degree. On Koblitz curves, the map
$\pi:(x,y)\mapsto(x^2,y^2)$ is also a curve endomorphism. Its action on the
selected cyclic subgroup can be represented by a known eigenvalue
$\lambda\bmod r$, allowing orbit-weighted relations.\cite{ggmp2020}

Over $\F_{p^k}$, a base defined by $x\in\F_p$ interacts with summation
polynomials through Weil restriction. Fixed-degree extension-field
algorithms and their symmetry reductions are established antecedents of
the implemented work. Their asymptotic opportunities do not imply
reasonable constants at the measured sizes, nor do they imply an
algorithm on $E(\F_p)$.\cite{gaudry2009,diem2011,fghr2014}

The implemented studies therefore cover prime fields, binary extensions,
and selected odd-characteristic extensions, notably degrees two through
six, together with protocol-specific extension arithmetic. This is not a
claim of a complete backend for every finite field, every characteristic
or every elliptic-curve model. Backend limits and unimplemented field
degrees remain separate from mathematical applicability.

\subsection{Isogenies and higher-dimensional transfers}

An isogeny defined over the working field is a homomorphism on rational
points. If its restriction to the prime subgroup is injective, then
$\varphi(Q)=[d]\varphi(G)$ preserves the logarithm. A destination solver
must still solve that transported instance more cheaply after charging
map discovery, construction, evaluation and recovery. A small degree,
shared $j$-invariant, common group order or common endomorphism field is
not by itself an executable transport certificate.

A cover or Jacobian route has more obligations: a concrete destination,
an explicit nonzero map, a proof or check of subgroup preservation, an
implementable destination solver, and the total source-to-answer cost.
Merely finding a low-genus curve with desirable arithmetic does not
establish that the source ECDLP can reach it. The portfolio's newer
identity and curve registries explicitly encode these distinctions.
\src{TRANSFER}\src{TEN}\src{VOLCANO}
